The proposed method for this project will be to ingest data from Fink, filtering the incoming data based on the scores that classification algorithms such as Early SN Ia, CATS and SuperNNova give to build an initial sample of SNe. Following on from that initial sample the method of \textcite{deckers_constraining_2022}. \textcite{deckers_constraining_2022} made cuts based on colour parameters derived from the application of the SALT2 model, where SNe Ia with \(-0.3 \leq c \leq 0.3\) are cut from the sample. Cuts are then made based on the stretch parameter \(-3.0\leq x_1\leq 3.0\), then a cut based on the peak magnitude in the g band of the observation \(-19.95\leq peak_g\leq-18.35\) and finally the first detection of the SNe Ia needs to be 14 days before the peak in the green bands luminosity was reached. \textcite{deckers_constraining_2022} argues that this is to remove peculiar SNe Ia from the data set and leaves only 'normal' SNe Ia in the dataset for analysis. 

It is unclear whether \textcite{deckers_constraining_2022} uses the following method for detection of EEx SNe Ia on the original starting sample of 127 SNe Ia or if it is on the sample of 115 SNe Ia post quality cuts. It is a reasonable assumption that the detection of EEx SNe Ia was done on the sample post quality cuts. 

The method to detect the EEx SNe Ia used in \textcite{deckers_constraining_2022} assess residuals in the light curves in flux space that are found by taking the difference between the best fit model light curve from the SALT2 model and the observed light curve. Specifically focusing on the span of time 6 days after the estimated explosion epoch which is determined from the best fit model. Although for this project, the SALT3 light curve model will be used due to its improved performance in the UV and blue wavelengths when compared to the SALT2 model. This is because these bumps, as mentioned in chapter \ref{sec: Supernova Ia physics and Progenitor Channels} subsection \ref{subsec: Early Flux Excess}, are universally near UV to blue throughout the excess.

Following the identification and classification of the EEx SNe Ia, the assigning of host galaxies will be done using the Directional Light Radius method, following this assignment the correlations of the EEx SNe Ia population in the sample with the properties of the host galaxies will be determined. These properties include but are not limited to the sSFR, stellar mass of the galaxy and the global colour of the galaxy. The colour of the EEx SNe Ia will also be a part of these correlative analyses to determine any correlations with host galaxy properties and the colour of the EEx SNe Ia. 